\section{Composition and exact expansions of original circuits}
\label{sec:peps_original_chronology_bounds}

The application of distributed compression to a dyadic protocol uses the
actual ordered composition of its changes. The following elementary identities
retain the original occurrences and local expansions. They are the chronological
comparisons used in Proposition~7.1 and in the assembly argument of
\cite{OpenAI2026PolynomialPEPS}; they do not construct the protocol.

\begin{definition}[Chronological composite]
    \label{def:peps_original_chronology_composite}
    \leanok
    \lean{TNLean.PEPS.PairEffect.OriginalCircuit.sequence}
    For original distributed circuits
    $W_t:\mathcal H_{A_t}\to\mathcal H_{A_{t+1}}$, define
    $W^{[0]}=1$ and $W^{[t+1]}=W_tW^{[t]}$.
\end{definition}

\begin{theorem}[Chronological composition]
    \label{thm:peps_original_chronology_composition}
    \leanok
    \uses{def:peps_original_chronology_composite}
    \lean{TNLean.PEPS.PairEffect.OriginalCircuit.eval_sequence_zero}
    \lean{TNLean.PEPS.PairEffect.OriginalCircuit.eval_sequence_succ}
    \lean{TNLean.PEPS.PairEffect.OriginalCircuit.isExpansionBounded_sequence}
    \lean{TNLean.PEPS.PairEffect.OriginalCircuit.nonprivateCount_sequence}
    \lean{TNLean.PEPS.PairEffect.OriginalCircuit.forall_participants_sequence}
    \lean{TNLean.PEPS.PairEffect.OriginalCircuit.isMonomialBounded_sequence}
    \lean{TNLean.PEPS.PairEffect.OriginalCircuit.participationCount_comp}
    \lean{TNLean.PEPS.PairEffect.OriginalCircuit.participationCount_sequence}
    Let $W_t:\mathcal H_{A_t}\to\mathcal H_{A_{t+1}}$, $0\le t<T$,
    be original distributed circuits.
    The composite acts by the ordered product: for
    $v\in\mathcal H_{A_0}$, $W^{[0]}v=v$ and
    $W^{[t+1]}v=W_t\bigl(W^{[t]}v\bigr)$.
    Write $N(W)$ for the number of nonprivate occurrences in $W$, and
    $n_p(W)$ for the number involving a fixed party $p$. Then
    \begin{align}
        N(W^{[T]})   & =\sum_{t<T}N(W_t),   \notag \\
        n_p(W^{[T]}) & =\sum_{t<T}n_p(W_t).\notag
    \end{align}
    Every property of the participating set that holds at every original
    occurrence also holds at every occurrence of the composite. If the original
    monomial lists have at most $K$ terms, at most $r$ pair effects per term,
    and coefficient sum at most $S$, the same bounds hold in the composite.
    Occurrences are retained separately even when their local operators agree.
\end{theorem}
\begin{proof}\leanok
    Induct on the number of stages. Appending the last circuit applies
    $W_t$ after $W^{[t]}$, which is the evaluation identity. Its occurrences are the disjoint union of the
    preceding occurrences and those of the last stage, with the same
    participating sets and monomial lists. Both counting identities and the
    stated bounds follow.
\end{proof}

\begin{definition}[Recorded local operators]
    \label{def:peps_original_recorded_gate_maps}
    \leanok
    \lean{TNLean.PEPS.PairEffect.OriginalCircuit.canonicalGateMaps}
    At every nonprivate occurrence of an original distributed circuit $W$,
    prescribe its recorded monomial sum as the local operator. Write
    $\widehat W$ for the operator obtained by evaluating the chronology of $W$
    with these prescribed local operators.
\end{definition}

\begin{theorem}[Exact local operator sums]
    \label{thm:peps_original_exact_gate_maps}
    \leanok
    \uses{def:peps_original_recorded_gate_maps}
    \lean{TNLean.PEPS.PairEffect.OriginalCircuit.evalWithGateMaps_canonicalGateMaps}
    \lean{TNLean.PEPS.PairEffect.OriginalCircuit.isGateApproximation_zero_canonicalGateMaps}
    For every original distributed circuit $W$, $\widehat W=W$, including the
    private operations. Every local approximation error is zero; the
    original contraction and pair-effect conditions are retained.
\end{theorem}
\begin{proof}\leanok
    At a nonprivate occurrence with monomial list $L$, the prescribed operator
    is $\sum_{m\in L}m$, which is the original gate, so $\widehat W=W$ for a
    single gate and its local error $\bigl\|\sum_{m\in L}m-\sum_{m\in L}m\bigr\|$
    vanishes. If $\widehat U=U$ and $\widehat V=V$, then
    $\widehat{VU}=\widehat V\,\widehat U=VU$ and
    $\widehat{1\otimes U}=1\otimes\widehat U=1\otimes U$; private operations
    and exchanges of registers are evaluated unchanged. Induction over the
    circuit gives $\widehat W=W$ and the zero local error assertion.
\end{proof}

\begin{theorem}[Original monomial and expansion bounds]
    \label{thm:peps_original_produced_resource_bounds}
    \leanok
    \lean{TNLean.PEPS.PairEffect.OriginalCircuit.monomialCount_produce_le_of_isMonomialBounded}
    \lean{TNLean.PEPS.PairEffect.OriginalCircuit.isExpansionBounded_mono}
    If every original monomial list has at most $K$ terms, then every
    corresponding produced occurrence has at most $K$ terms. If an original
    circuit has pair-effect bound $r$ and coefficient-sum bound $S$, then
    bounds $r'\ge r$ and $S'\ge S$ also bound that circuit.
\end{theorem}
\begin{proof}\leanok
    The original lists are unchanged at each occurrence. Their length bound
    gives the first assertion. For the second, enlarge the pair-effect bound
    in each monomial and the coefficient-sum bound by transitivity.
    Composition and spectator identities preserve these assertions.
\end{proof}
